%%%%%%%%%%%%%%%%%%%%%%%%%%%%%%%%%%%%%%%%%%%%%%%%%%%%%%%%%%%%%%%%%%%%%%%%%%%%%%
% Fiche d'évaluation AMC — Cellule robotisée
%
% TP noté d'Automatisme pour la robotique — BUT GEII
% Mode standalone (génération par CSV, sans extraction automatique du TP)
%
% Structure :
%   - 14 critères répartis en 5 groupes (20 pts au total)
%   - Groupe robot       : Prise/Dépose Robot              (8 pts)
%   - Groupe pont_config : Configuration réseau             (3 pts)
%   - Groupe pont_prog   : Programmation communication      (5 pts)
%   - Groupe pont_mouvt  : Mouvements du chariot            (2 pts)
%   - Groupe pont_integ  : Intégration programme principal  (2 pts)
%
% Compilation :
%   pdflatex fiche_evaluation.tex
%   (nécessite automultiplechoice, csvsimple, UPSTI_Document, UPSTI_AMC)
%%%%%%%%%%%%%%%%%%%%%%%%%%%%%%%%%%%%%%%%%%%%%%%%%%%%%%%%%%%%%%%%%%%%%%%%%%%%%%

\documentclass[QCM, noCustomPackages]{UPSTI_Document}

% ============================================================
% PROFIL ET CONFIGURATION INSTITUTIONNELLE
% ============================================================
\UPSTIloadProfile{BUT_GEII}

% ============================================================
% MÉTADONNÉES DU DOCUMENT
% ============================================================
\discipline[Autom - Robotique]{Automatisme pour la robotique}

% Titre affiché dans la grande zone centrale de l'en-tête UPSTI
\newcommand{\UPSTItitreEnTete}{Évaluation pratique}

% Titre du document (affiché sous l'en-tête en mode QCM)
\newcommand{\UPSTItitre}{Cellule robotisée -- Évaluation pratique}

% Numéro/type affiché à droite de l'en-tête (colonne type de document)
\newcommand{\UPSTInumero}{Noté}

% Numéro de séquence (commande locale à ce projet, non UPSTI)
\newcommand{\sequence}{99}

% Version E = version étudiante (cases en bleu clair, pas de correction)
\documentVersion{E}

% Valeur par défaut (sera surchargée dans chaque \exemplaire par \id{})
\newcommand{\UPSTInumeroVersion}{--}

% Message affiché dans la zone basse de l'en-tête
% (sera surchargé dans chaque \exemplaire par le nom de l'étudiant)
\newcommand{\UPSTImessage}{}

% ============================================================
% DÉBUT DU DOCUMENT
% ============================================================
\begin{document}

% Graine aléatoire AMC : fixe l'ordre de tirage des questions
% (valeur choisie arbitrairement, à conserver pour reproductibilité)
\AMCrandomseed{202599}

% Mode de tirage des groupes : sans remise (ordre fixe, tous les critères)
\setdefaultgroupmode{withoutreplacement}

% ============================================================
% BANQUE DE CRITÈRES
% ============================================================
% Syntaxe : \UPSTIcritere[pts]{id}{description}
%   - pts         : points accordés pour ce critère (défaut : 1)
%   - id          : identifiant unique (utilisé par AMC pour la correction)
%   - description : texte affiché sur la fiche étudiante
%
% Chaque \UPSTIcritere est encapsulé dans un \element{groupe}{...}
% \restituegroupe{groupe} pioche dans cette banque lors de la génération.
% ============================================================

% ------------------------------------------------------------
% Groupe : robot — Partie 1 - Prise/Dépose Robot (12 pts)
% ------------------------------------------------------------

\element{robot}{%
  \UPSTIcritere[2]{var-robot}{%
    Variables de position déclarées avec les bons types et valeurs initiales.%
  }%
}

\element{robot}{%
  \UPSTIcritere[2]{shift-post}{%
    Bloc \texttt{VAL\_ShiftPoint} instancié et appelé au bon endroit
    dans le GRAFCET coopératif.%
  }%
}

\element{robot}{%
  \UPSTIcritere[3]{shift-calcul}{%
    Branchements de \texttt{VAL\_ShiftPoint} corrects :
    entrées et sortie du bloc correctement câblées.%
  }%
}

\element{robot}{%
  \UPSTIcritere[2]{mouv-approche}{%
    Commande de mouvement vers le point d'approche correctement conditionnée.%
  }%
}

\element{robot}{%
  \UPSTIcritere[3]{grafcet-complet}{%
    GRAFCET \texttt{F1\_EVAL\_PriseDepose} complet et comportement conforme
    au cahier des charges.%
  }%
}

% ------------------------------------------------------------
% Groupe : pont — Partie 2 - Pont roulant (6 pts, 3 critères)
% ------------------------------------------------------------

\element{pont}{%
  \UPSTIcritere[2]{pont-config}{%
    Configuration réseau correcte : module ajouté dans le DTM Navigator,
    tailles E/S renseignées, adresse IP sur le bon coupleur NOC.%
  }%
}

\element{pont}{%
  \UPSTIcritere[2]{pont-comm}{%
    Programmation de la communication : variables déclarées, blocs de
    lecture/écriture câblés et positionnés correctement dans le cycle API.%
  }%
}

\element{pont}{%
  \UPSTIcritere[2]{pont-fonct}{%
    Fonctionnement du chariot : bloc \texttt{EVAL\_MoveCart} correctement
    appelé et séquence principale robot\,+\,pont opérationnelle.%
  }%
}

% ------------------------------------------------------------
% Groupe : qualite — Qualité du code (2 pts)
% ------------------------------------------------------------

\element{qualite}{%
  \UPSTIcritere[2]{qualite-code}{%
    Qualité générale du code : nommage cohérent, structure lisible,
    absence de code mort ou de câblages inutiles.%
  }%
}

% ============================================================
% GÉNÉRATION DES COPIES NOMINATIVES
% ============================================================
% Pour chaque ligne du CSV (colonnes : id, prenom, nom) :
%   - \renewcommand{\UPSTImessage}{...} : affiche le nom dans l'en-tête
%   - \renewcommand{\UPSTInumeroVersion}{...} : identifiant étudiant (AMC)
%   - \UPSTIbuildPage : construit l'en-tête UPSTI
%   - \restituegroupe{...} : tire tous les critères du groupe (sans remise)
%   - \AMCcleardoublepage : saut de page entre deux copies
% ============================================================

\csvreader[head to column names]{etudiants.csv}{}{%
  \exemplaire{1}{%

    % --- Réinitialisation du compteur de critères ---
    \UPSTIresetCritere
    % --- Personnalisation de l'en-tête pour cet étudiant ---
    \renewcommand{\UPSTImessage}{\champnom{\nom\ \prenom}}%
    \renewcommand{\UPSTInumeroVersion}{\id{}}%
    \UPSTIbuildPage
    \AMCassociation{\id}


    % --------------------------------------------------------
    % Partie 1 — Robot : Prise/Dépose (12 pts)
    % --------------------------------------------------------
    \section*{Partie~1 --- Robot : Prise/Dépose
              \hfill \normalfont\small /12~pts}

    \restituegroupe{robot}

    \smallskip
    \noindent\textit{Commentaires :}\\[2pt]
    \UPSTIquadrillage{9}

    \bigskip

    % --------------------------------------------------------
    % Partie 2 — Pont roulant et traitement de surface (6 pts)
    % --------------------------------------------------------
    \section*{Partie~2 --- Pont roulant et traitement de surface
              \hfill \normalfont\small /6~pts}

    \restituegroupe{pont}

    \smallskip
    \noindent\textit{Commentaires :}\\[2pt]
    \UPSTIquadrillage{3}

    \bigskip

    % --------------------------------------------------------
    % Qualité du code (2 pts)
    % --------------------------------------------------------
    \section*{Qualité du code
              \hfill \normalfont\small /2~pts}

    \restituegroupe{qualite}

    \smallskip
    \noindent\textit{Commentaires :}\\[2pt]
    \UPSTIquadrillage{2}

    % Saut de page entre deux copies (fin de recto ou verso)
    \AMCcleardoublepage
  }%
}

\end{document}
